\color{unidarkgreen}{\subsection{Общая логика}}
\color{black}

\begin{enumerate}

\item
Общая логики формирует общие выходные сигналы:
\begin{itemize}
\item \textit{«Терминал выведен»};
\item \textit{«Неиспр. терминала»};
\item \textit{«Неиспр. ОТ терм»};
\item \textit{«БИ выведены»};
\item \textit{«Несоответ. ПУ»};
\item \textit{«Неиспр. ОТ терм»};
\item \textit{«Вых. цепи разобраны»}.
\end{itemize}

\item
Алгоритм работы общей логики приведен на рисунке \ref{pic:GL}. Уставки общей логики приведены в таблице \ref{app_set:obshlog}.

\medskip
\begin{figure}[H]
\centering
\includegraphics[width=0.9\textwidth,height=0.9\textheight,keepaspectratio]{\fbpath/05.01 ДЗО/90 Общая логика/img/GL.pdf}
\captionof{figure}{Функциональная схема общей логики}\label{pic:GL}
\end{figure}


\end{enumerate}